\documentclass[12pt]{article}
\usepackage[utf8]{inputenc}
\usepackage[T1]{fontenc}
\usepackage{hyperref}
\title{Kanjari Language Resources for K-12 Curriculum Development}
\author{Compiled by Claude Sonnet 5.5 (claude-sonnet-5-5) for Andr\'as Kornai}
\date{}
\begin{document}
\maketitle

\section{Introduction}
Kanjari (ISO 639-3 kft; Glottocode kanj1259) is an Indo-Aryan language associated with the Kanjar community of India and Pakistan, with about 210,000 speakers (2011). Classification varies between sources. Documentation is very limited.

\section{Primary Educational Resources}
None found.

\section{Linguistic and Reference Materials}
\begin{itemize}
\item Wikipedia, Kanjari language: \url{https://en.wikipedia.org/wiki/Kanjari_language}
\item Wikipedia, Kanjar (people): \url{https://en.wikipedia.org/wiki/Kanjar}
\item Ethnographic overview, ``Kanjar -- Orientation'': \url{https://www.everyculture.com/South-Asia/Kanjar-Orientation.html}
\end{itemize}
Grierson's Linguistic Survey of India is mentioned as a source, but no URL appeared in the results. No modern grammar or dictionary was found.

\section{Religious and Community Resources}
None found.

\section{Digital and Audio Resources}
None found.

\section{Curriculum Development Context}
The language is reported as endangered by one source. There is no material suitable for classroom use, so any curriculum would require original documentation with the community.
\end{document}